\quellseite{33}
\chapter{A. Religious Instruction in the Early Church}

*J. W. F. Hoefling, \emph{Das Sakrament der Taufe nebst den andern damit zusammenhaengenden Akten der Initiation} i, 1846. Th. Harnack, \emph{Der christliche Gemeindegottesdienst im apostolischen und alt-katholischen Zeitalter}, 1854. *G. v. Zezschwitz, \emph{Katechetik} i and ii, 1. J. Mayer, \emph{Geschichte des Katechumenats und der Katechese in den ersten sechs Jahrhunderten}, 1868. Fr. Probst, \emph{Lehre vom Gebet in den drei ersten christlichen Jahrhunderten}, 1871. Fr. Probst, \emph{Die kirchliche Disziplin in den drei ersten Jahrhunderten}, 1873. Fr. Probst, \emph{Katechese vom Anfang des vierten bis zum Ende des sechsten Jahrhunderts}, 1884. Fr. Probst, \emph{Geschichte der katholischen Katechese}, 1886. *H. J. Holtzmann, \emph{Die Katechese der alten Kirche}, 1892. E. Sachsse, \emph{Katechetik}, pp. 4–115. *F. Wiegand, \emph{Die Stellung des apostolischen Symbolums im kirchlichen Leben des Mittelalters}, i: Symbol und \emph{Katechumenat}, 1899. F. Cohrs, \emph{Katechumenat}, PRE³, 1901. W. Moeller-Schubert, \emph{Lehrbuch der Kirchengeschichte} i, ² 1902. J. H. Kurtz-Bonwetsch, \emph{Lehrbuch der Kirchengeschichte} i, ¹⁴ 1906. G. Hodgson, \emph{Primitive Christian Education}, 1906. J. C. Ayer, \emph{Catechumenal and Catechetical Schools} (Monroe’s Cyclopedia, vol. 1). A. F. Leach, \emph{Bishop’s Schools and Cathedral Schools} (ibid., vol. 1). G. Krueger, \emph{Handbuch der Kirchengeschichte} i, 1911. E. Chr. Achelis, \emph{Prakt. Theol.}, ii, ³ 1911. H. Achelis, \emph{Das Christentum in den ersten drei Jahrhunderten}, 1912. J. Steinbeck, \emph{Katechetik}, pp. 1–9. Migne \emph{Patrologiae Cursus completus} (MSG = Greek Series, MSL = Latin Series). \emph{Corpus Scriptorum Ecclesiasticorum Latinorum}, Vindobonae (CSEL). \emph{Die griechischen christlichen Schriftsteller der ersten drei Jahrhunderte}, Berlin (GBSchr). \emph{Bibliothek der Kirchenvaeter}, Kempten, ed. by Thalhofer (KB¹). \emph{Bibliothek der Kirchenvaeter}, Kempten, ed. by Bardenhewer, Scherman, and Weyman (KB²). \emph{The Ante-Nicene Fathers} (ANF). \emph{Nicene and Post-Nicene Fathers}, ed. by Schaff, first and second series (NPNF¹, NPNF²). A. Harnack, \emph{Geschichte der alt-christlichen Literatur bis Eusebius}, 1893–1904. Bardenhewer, \emph{Geschichte der altkirchlichen Literatur}, 1902 ff. *H. Jordan, \emph{Geschichte der alt-christlichen Literatur}, 1911.

\quellseite{34}
\section[§ 4. The Aim in the Early Church]{§ 4. The Aim of Religious Instruction in the Early Church}

G. Thomasius-Bonwetsch, \emph{Christliche Dogmengeschichte} ¹, 1886. M. v. Nathusius, \emph{Handbuch des kirchlichen Unterrichts nach Ziel, Inhalt und Form} ¹, 1903. F. Loofs, \emph{Leitfaden zum Studium der Dogmengeschichte}, ⁴ 1906. A. Harnack, \emph{Mission und Ausbreitung des Christentums in den drei ersten Jahrhunderten}, 1906, ⁴ 1924. F. Kattenbusch, Taufe, PRE ³, 1907. *R. Seeberg, \emph{Lehrbuch der Dogmengeschichte} ¹, ³ 1922. A. Harnack, \emph{Lehrbuch der Dogmengeschichte} ¹, ⁴ 1909. A. v. Stromberg, \emph{Studien zur Theorie und Praxis der Taufe in der christlichen Kirche der ersten zwei Jahrhunderte}, 1913.

The early Church was, to a degree never excelled, a missionary organization, and was, therefore, primarily interested in the induction of adults into membership. In conformity with the command of Matt. 28:19 and with apostolic usage, reception into the Church took place through the act of baptism. This sacred act was highly evaluated in direct proportion as the Church vitally recognized the unique character of Christianity. It was regarded as the agency that worked forgiveness of sins and terminated the old life, as the channel for the Holy Spirit and the beginning of man’s renewal. It was viewed as the door through which access was afforded to all the blessings and gifts in possession of the Church of God on earth. And when this consciousness had largely begun to fade, the praise of the teachers of the Church for baptism as an agency of salvation became particularly marked.

However, the Early Church made the saving effect of the baptismal sacrament contingent on personal faith, and at the same time viewed it as the strongest incentive to a godly life. Thus the urgent necessity arose to impart religious instruction and training to all prospective members
\quellseite{35}
with a view to a blessed reception of baptism. The function of the missionary had to give place to that of the catechist in order to prepare those for baptism who were willing to join the congregation of Christian believers. In this way specific forms of a catechumenate, or a baptismal discipline, gradually developed which, notwithstanding minor local differences, universally possessed the same fundamental features. The earliest traces of the catechumenate are found in the apostolic era; it attained its highest perfection in the fourth century, but soon after, already in the fifth century, it began to decay.

\quellseite{36}
\section[§ 5. Religious Instruction During the Apostolic Age]{§ 5. Religious Instruction During the Apostolic and Post-Apostolic Age}

*G. v. Zezschwitz i, § 11. 12; ii, 1, § 16. 17. 19. 20. 26. 31. *Th. Zahn, \emph{Glaubensregel und Taufbekenntnis in der alten Kirche}, 1881 (in \emph{Skizzen aus dem Leben der alten Kirche}, ³1908: compare his article \emph{Glaubensregel} in PRE³, 1899). W. Bacher, \emph{Die Agada der Tannaiten} i 1884, ²1903, ii 1900. R. Harris, \emph{The Teaching of the Apostles and the Sibyll. Books}, 1885. C. Taylor, \emph{The Teachings of the Twelve Apostles with Illustrations from the Talmud}, 1886. W. Bacher, \emph{Die Agada der palaest. Amoraeer}, 1892–1899. A. Harnack, \emph{Das apostolische Glaubensbekenntnis}, 1892, ²1896. J. H. Holtzmann, pp. 62–66. *Th. Zahn, \emph{Das apostolische Symbolum}, 1893. *F. Kattenbusch, \emph{Das apostolische Symbolum}, i 1894, ii 1900. A. Harnack, \emph{Apostolisches Symbol} in PRE³, 1896. A. Harnack, \emph{Die Apostellehre und die juedischen beiden Wege}, ²1896 (compare \emph{Apostellehre} in PRE³, 1896). A. Hahn, \emph{Bibliothek der Symbole und Glaubensregeln der alten Kirche}, ³1897. *J. Kunze, \emph{Glaubensregel, Heilige Schrift und Taufbekenntnis}, 1899. C. Weizsaecker, \emph{Das apostolische Zeitalter der christlichen Kirche}, ²1902. A. C. McGiffert, The Apostles' Creed, 1902. W. Heitmueller, \emph{Im Namen Jesu}, eine sprach- und relig.-geschichtliche Untersuchung, spez., zur altchristlichen Taufe, 1903. *A. Seeberg, \emph{Der Katechismus der Urchristenheit}, 1903. *A. Seeberg, \emph{Das Evangelium Christi}, 1905. F. W. Rendorff, \emph{Die Taufe im Urchristentum}, 1905. R. Knopf, \emph{Das nachapostolische Zeitalter}, 1905. F. v. Dobschuetz, Proselyten in PRE, 1905. A. Harnack, \emph{Mission und Ausbreitung des Christentums}, 1906, ⁴1924. *A. Seeberg, \emph{Die beiden Wege und das Aposteldekret}, 1906. *A. Seeberg, \emph{Die Didache des Judentums und der Urchristenheit}, 1908. *R. Seeberg, \emph{Dogmengeschichte} i, § 6.8. A. Harnack, \emph{Dogmengeschichte} i, pp. 172 ff.; 225 ff. M. Reu, \emph{Der katechetische Unterricht in der apostolischen Zeit} (Kirchliche Zeitschrift), 1909. F. Schuerer, \emph{Geschichte des juedischen Volkes im Zeitalter Jesu} iii, ⁴1909. G. Klein, \emph{Der aelteste christliche Katechismus und die juedische Propagandaliteratur}, 1909. W. Brandt, \emph{Die juedischen Baptismen}, 1910. J. Kunze, \emph{Das apostolische Glaubensbekenntnis und das Neue Testament}, 1911. *J. Behm, \emph{Handauflegung im Urchristentum}, 1911. W. Heitmueller, \emph{Taufe und Abendmahl im Urchristentum}, 1911. O. Moe, \emph{Paulus und}
\quellseite{37}
\emph{die evangelische Geschichte}, 1912. A. Seeberg, \emph{Die Taufe im Neuen Testament}, ² 1912. H. Achelis, \emph{Das Christentum} etc., ¹, pp. 121 ff. A. v. Stromberg, \emph{Studien}, etc., 1913. H. Lietzmann, \emph{Symbole der alten Kirche}, 1914. J. Haussleiter, \emph{Trinitarischer Glaube und Christusbekenntnis in der alten Kirche}, 1920. *Feine, \emph{Die Gestalt des apost. Glaubensbekenntnisses in der Zeit des N. T.}, 1925. *W. Foerster, \emph{Herr ist Jesus. Herkunft und Bedeutung d. urchristlichen Kyrios-Bekenntnisses}, 1925. Sources: \emph{Patrum apostolicorum opera}, rec. Oscar de Gebhardt, Ad. Harnack et Theod. Zahn, editio minor ⁵1906; German by I. Chr. Mayer, 1869 (KB¹); English in ANF, vol. i. \emph{The Didache or Teachings of the Twelve Apostles}: \emph{Didache mit kritischem Apparat} ed. by H. Lietzmann, ³1912; English in ANF, vol. vii, pp. 369 ff.; Greek and German ed. by Harnack, 1896. \emph{Writings of the Apologists}: E. J. Goodspeed, \emph{Die Apologeten, die aeltesten Texte mit kurzen Einleitungen}, 1914; English in ANF, vol. i and ix; German by Rauscher, Julius, a. o. in KB², vol. i and ii. \emph{The Apologies of Justin}: \emph{Justins des Maertyrers Apologien}, ed. G. Krueger, ⁴1904; English in ANF, vol. i; German by Rauscher in KB², vol. i.

At the beginning of his ministry Jesus called the Twelve to become fishers of men, to induce others to become disciples of Jesus (Matt. 4: 19; Luke 5: 10). Already during his sojourn in the flesh they began to exercise their missionary vocation (John 4: 12; Matt. 9: 36—10: 15; Mark 6: 7–13; Luke 9: 1–6). Despite their failure to meet the test during his passion, the risen Christ reinstated his disciples in the apostolate (Matt. 28: 19; Luke 24: 47; John 20: 21–23; 21: 15–23; Acts 1: 8) and assigned to them in unambiguous words the whole world, Jews and Gentiles, as the object of their missionary activity, thereby dispelling whatever remnant of uncertainty in this matter may still have lingered in their minds. To make them sure of their ground, the Saviour designated the means of discipling the nations, viz., the testimony concerning him, baptism into the name of the triune God, and the inculcation of all his commands.

The Lord, indeed, had nowhere explicitly declared instruction to be a necessary part of the preparation for baptism. The words sometimes adduced as proving that the Lord
\quellseite{38}
enjoined such instruction are not in point: the verbs \emph{κηρύσσειν} (Luke 24:27) and \emph{μαρτυρεῖν} (Acts 1:8) denote missionary proclamation; the \emph{βόσκειν} and the \emph{ποιμαίνειν} of the lambs of Jesus (John 21:15) presuppose membership in the congregation and, accordingly, also baptism whereby it is effected; and the order of sentences in Matt. 28:19 f. also shows that the work of “teaching them to observe all things whatsoever I have commanded you,” follows upon baptism. Especially the last passage cannot have failed to remind the apostles of the vast amount of disciplinary toil which their Master was forced to lavish upon them even long after they had attained to faith; and, undoubtedly, these words suggested to them the duty of continuously training and instructing those who by baptism had been inducted into the Christian fold so that their walk of life might ever more fully conform to their confession. But a specific command to impart pre-\emph{baptismal instruction} was not needed at all; for if it is true that baptism removes man into communion with Father, Son, and Spirit, then pre-\emph{baptismal instruction} regarding these factors of salvation simply cannot be escaped, and if the baptised persons were expected to lead holy lives, then they must have been impressed with the urgent necessity of a profound change in morals and they must have been given at least rudimentary information about the \emph{ἐντολαί} of Christ, before the Sacrament was administered.

It was likewise obvious that \emph{baptismal instruction} as imparted to \emph{Gentiles} must differ from that imparted to \emph{Jews}. In the case of the latter there probably was no specific catechetical activity at all, for when the missionary’s sermon had convinced the Jewish hearer that Jesus is the promised Christ sent by the Father, and when he had expressed his readiness to accept him as his Saviour and King, all the conditions for a blessed reception of baptism were fulfilled, particularly since the ethical standard required by Christ was
\quellseite{39}
not at all contrary to that of the Old Covenant. Thus we observe that Peter’s catechetical work on the day of Pentecost was limited to what Luke reports in Acts 2:38–40, and Philip’s way of treating Queen Candace’s treasurer was not radically different because as a proselyte he had been “precatechised by the prophets.” The baptismal instruction of Gentiles must needs be much more comprehensive, particularly when there had been no acquaintance with the Septuagint and Judaism. Here both factors called for due recognition: instruction concerning salvation as wrought by God through Christ and appropriated in the Spirit, and instruction concerning the moral life required of a disciple of Christ.

Should the apostles have had any doubt concerning the proper arrangement of the baptismal instruction for Gentiles, the proselyting method in vogue among the Jews readily suggested itself as a sort of model. For decades Jewish missionaries had covered the mediterranean world (Matt. 23:15); they had developed a highly successful propaganda and created a genuine catechumenate for proselytes, which embraced instruction in faith and morals and concluded with a baptism. It is easily explained that the apostles, who also in other respects, in worship and private life, clung to the historic forms of Judaism, adopted this \emph{proselyte catechumenate} as far as their convictions permitted.

The chief features of the \emph{proselyte catechumenate} were as follows: When anyone had declared his willingness to become a member of the covenant people, he first received instruction in regard to the God that had created the world and wondrously guided Israel. Thereby he took upon himself “the yoke of the Kingdom of God.” This was followed by instruction in regard to right morals based chiefly on Lev. 18 and 19; its substance may still be found in the \emph{“Two Ways”} of the \emph{Didache}. Accepting this ethical standard, the proselyte submitted to “the yoke of the commandments.” Probably an eschatological element, treating of final judgment and Israel’s ultimate glory was also included in this instruction. The preparation was followed
\quellseite{40}
by the tripartite act of reception: 1. Circumcision (followed by a feast); 2. the \emph{thebilah}, or baptism by immersion; 3. a sacrifice, or, in the dispersion, fasting and alms.

A careful examination of the New Testament actually discloses the fact that before baptism a certain amount of instructive material was delivered to every Christian which the apostles utilized as the basis for further instruction or admonition. To this common foundation the author of Hebrews refers when he speaks of the “first word of Christ” through which the foundation of the readers’ faith has been laid (6:1), for this can only signify the instruction by means of which the Saviour, operating through the apostles and their co-laborers had begun his work in them. A considerable number of passages confirms the existence of such baptismal instruction, esp. 1 Thess. 4:1; 2 Thess. 3:6; 1 Cor. 4:17 (\emph{αἱ ὁδοί μου ἐν Χριστῷ}); Rom. 6:17 (analyze as follows: \emph{εἰς τὸν τύπον (τῆς) διδαχῆς, ὃν παρεδόθητε = ὃς παρεδόθη ὑμῖν}); Rom. 16:17; Eph. 4:20 (\emph{διδάσκεσθαι καὶ μανθάνειν τὸν Χριστόν}). The terms \emph{παράδοσις}, \emph{παραδιδόναι} and the corresponding terms \emph{παραλαμβάνειν} and \emph{μανθάνειν} become particularly significant if we bear in mind that among the Jews \emph{παράδοσις} had become the technical term for the oral transmission of the law (Mark 7:3). While the above passages refer to instruction in morals, 1 Cor. 15:3–5; 1 Tim. 6:11 ff.; 2 Tim. 3:10—4:3 (Compare also Rom. 6:4; Col. 2:11–13; 1 Peter 3:18—22) refer to instruction in the faith; and Heb. 6:2 (cf. \emph{Didache} 16) shows that the eschatological element was not lacking.

The literature extant enables us to ascertain rather definitely the contents of this body of evangelical truths and \emph{practical precepts}. Let 1 Thess. 4:3–8 serve as an illustration. Here we find that the apostle warns the Thessalonians against adultery, covetousness, and avarice; but he does so immediately after reminding them of the fact that he had previously given them moral precepts (v. 1: καθώς παρελάβετε παρ’ ἡμῶν etc.;
\quellseite{41}
v. 2: \emph{οἴδατε, τίνας παραγγελίας ἐδώκαμεν ὑμῖν}). This forces us to conclude that the warnings here given formed part of the body of moral precepts. But we meet the same sins time and again in the so-called catalogues of vices (Gal. 5:19 ff.; 1 Cor. 6:9 f.; Eph. 5:3–5; Col. 3:5 ff.; Rom. 1:29 ff.; 1 Cor. 5:10 ff.; 1 Tim. 1:9 ff.; Rev. 21:8; 22:15), and everywhere they appear with rather definite regularity, though not with absolute uniformity. A similar definiteness may be observed wherever the so-called catalogues of virtues (e.g., Gal. 5:22 ff.; Eph. 5:9) and the tables of domestic duties (e.g., Eph. 5:22—6:9; Col. 3:18 ff.; 1 Peter 3:1–7) appear. This suggests very strongly that such passages represent portions of the groups of moral precepts. According to the varying needs of the congregations addressed, the apostles would emphasize their admonitions by reminding them of the body of moral precepts delivered to them and by actually enumerating some of the sins and virtues contained therein, just as the preacher of today will occasionally lend emphasis to an admonition by quoting part of the Catechism. Fortunately, we possess, moreover, since 1883 in the “Two Ways” of the \emph{Didache} (ch. 1–6) an illustration of the moral formulas then in vogue and a cogent proof for the correctness of the above conclusions. At the same time, the Christians’ far-reaching dependence upon the Jewish proselyte catechumenate becomes evident, for the “Two Ways” represent a Christian adaptation of an originally Jewish catechism for proselytes.—It would seem that already in the days of Paul this ethical material was designated by definite terms; note, besides the expressions “way” and “doctrine” (Rom. 6:17; 16:17; 1 Cor. 4:17; Acts 9:2; both terms are used in a wider and narrower sense), especially the phrases “sound doctrine” (1 Tim. 1:10; 2 Tim. 4:3; Tit. 2:1) and “the good doctrine” (\emph{διδασκαλία}, 1 Tim. 4:6; cf. 5:17). At a later time (cf. Justin, Apologia i 14–17) the “doctrines of Christ” (διδάγματα τοῦ
\quellseite{42}
\emph{Χριστοῦ}) were frequently positively identified with his moral precepts (Matt. 28:19 and the Sermon on the Mount).

The contents of \emph{doctrinal instruction} may be deduced from the creedal statement connected with baptism. It is true that this formula is nowhere in the N. T. recorded in its entirety, and in the earliest times it can hardly have been perfectly uniform and unchangeable; nevertheless, such passages as 2 Tim. 3:10 ff.; 1 Tim. 6:13 ff.; 1 Cor. 15:3:4 permit of reasonably certain inferences. Accordingly, it must have treated first of God and then of Christ, all stress falling upon the christological features. It appears that in this formula God was confessed as the Creator of every living thing, and Christ as sent by Him (?), as coming from the seed of David, making his confession before (or in the time of) Pontius Pilate, having died for our sins according to the Scriptures, buried, having risen on the third day according to the Scriptures, appearing to Cephas and afterward to the Twelve, ascended to heaven (?), having under him the angels, principalities, and powers, and returning to judge the quick and the dead. This survey shows that instruction concerning God and Christ was an integral part of doctrinal instruction. The doctrine of God probably differed but little from that found in the Jewish catechumenate for proselytes. What was stressed, however, was the fact that God fulfilled his promise and sent His Son into the world for its salvation. The repeated usage of the phrase “according to the Scriptures” (1 Cor. 15: 3 ff.) indicates that the argument based upon prophecy and fulfillment which later became one of the favorite proofs for Christianity (cf. e. g., Justin \emph{Apol.} i 32 etc., and the \emph{Epideixis} of Irenaeus), played an important rôle in the baptismal instruction of the earliest times. In this connection the catechumen doubtless was introduced to the chief facts of the life of Jesus and to its significance as part and parcel of the work of salvation (1 Cor. 15: 3: ὑπὲρ τῶν ἁμαρτιῶν ἡμῶν).
\quellseite{43}
Here are the first traces of \emph{Biblical History} as catechetical material. In discussing the significance of Christ, however, strong emphasis must have been placed upon the truth that he now rules as sovereign over everything in heaven and on earth, and that he is to be worshipped as the One true “Lord,” in whom we must be saved. Thus only can we account for phrases which sound almost like formulas (1 Pet. 3:22; Eph. 1:20 f.; Col. 2:10; Rom. 8:34), and for the terminology used with reference to the confession of Christ as Lord in Rom. 10:9; Phil. 2:11; 1 Cor. 12:3; cf. Acts 2:36.—It appears that also this doctrinal material was gradually given definite designations; besides \emph{διδαχή} and \emph{ὁδός} which occur also in a wider sense, \emph{παραθήκη} (1 Tim. 6:20; 2 Tim. 1:12, 14) and \emph{πίστις} (Gal. 1:23; 1 Tim. 6:21; Jude 20) are so used, perhaps also \emph{εὐαγγέλιον} (Rom. 2:16; 2 Tim. 2:8).

The inferences thus far drawn are corroborated in every way by Heb. 6:2 no matter what reading (\emph{διδαχῆς} or \emph{διδαχήν}) or what construction is adopted. Both elements, moral and doctrinal instruction, are here mentioned side by side; “turning the heart away from dead works” denotes the former, and “turning the heart to God” the latter. The eschatological element is divorced from the others, very likely on account of the peculiar circumstances of the addressees. In one respect, however, this passage advances our investigation. So far no explicit statement referring to \emph{instruction concerning baptism} had been noted, but here the author specifically includes in prebaptismal instruction “the doctrine of baptisms and of the imposition of hands.” That the plural form “baptisms” is used, need not surprise us; the Christian teacher could not very well explain the peculiar nature of baptism without referring to the numerous washings practiced at that time, the baptism of John, the baptism of proselytes, the lustrations connected with various cults, etc. The “imposition of hands” can only refer to that applied at baptism. It is possible
\quellseite{44}
that the author of Hebrews regarded remission of sins as the special blessing of baptism, and the impartation of the Spirit as that of the imposition of hands. Others combined both effects with the one act of baptism (Acts 8:17; 2:38. So much is clear, however, that in discussing baptism the early Christians must have treated of remission of sins and the impartation of the Spirit. This fact is borne out by some passages in the New Testament which otherwise would be rather unintelligible (cf. Gal. 3:1–14 where the reception of the Holy Spirit is named as the specifically new blessing of Christianity; Eph. 1:13 where attaining to faith and receiving the seal of the Spirit are so closely connected; 1 John 2:20. 27 where the anointing from the Holy Spirit is practically identical with baptism). Quite likely anointing, sealing, and baptism were synonymous for the early Christians. Compare also Acts 19:2 ff.; Gal. 4:4–6; Rom. 8:15, and do not overlook the close connection of baptism with adoption and sonship.

It is no less certain that the \emph{words of institution of the Lord’s Supper} were imparted to the catechumens. We would draw such a conclusion from the order of parts in the \emph{Didache} (1–6 instruction, 7 baptism, 8 Lord’s Prayer, 9 Lord’s Supper), and also from the dependence of the Christians on the Jewish catechumenate (circumcision was followed by a festive meal). But we need not rely on inferences; Justin (\emph{Apol.} i 65) testifies that the newly baptized were admitted to Holy Communion, we know that in even earlier times the Lord’s supper was celebrated throughout the Christian church, and Paul explicitly declares that he transmitted the words of institution to the Christians at Corinth (1 Cor. 11:23). We are not so certain as to the time when these words were imparted, whether before or after baptism.

Finally, it appears that either shortly before, or immediately after baptism the catechumens were made acquainted with the Lord’s Prayer. The order of parts in the \emph{Didache}
\quellseite{45}
and later usage suggest this conclusion, to say the least, and passages like Rom. 8:15 and Gal. 4:6 corroborate it. Here Paul evidently reminds his readers of a well known prayer which they offer in their common services and which begins with the word Father. This supposition also serves to explain the juxtaposition of the Aramaic “Abba” and the Greek \emph{ὁ πατήρ}: along with Aramaic terms like Amen, Halle-lujah, Hosannah, Maranatha, also “Abba” was used in the common service of Greek speaking Christians; adding the Greek equivalent they evidently used this word to designate the Lord’s Prayer and to begin it. All the facts bearing upon the matter favor this construction: having become children of God through baptism and having received the Holy Spirit, they offer the prayer of the children of God.

\emph{Thus it is clear that the elements of instruction employed by the Church today, are rooted in apostolic times.}

Baptismal instruction was followed by baptism. It is not difficult to reconstruct the \emph{baptismal rite}. A confession of sin introduced the holy act (cf. Acts 2:38 and John’s baptism), perhaps also an act of renunciation (Tit. 2:12, \emph{ἀρνησάμενοι}) and a confession of faith (1 Tim. 6:12; Heb. 10:22 f.). In case the latter was pronounced by the baptizer, it is probable that the catechumen signified his assent with the words, “Lord is Jesus” (1 Cor. 12:3; Rom. 10:9). Concerning the administration of baptism itself the \emph{Didache} writes (ch. 7): “Baptize into (\emph{εἰς}) the name of the Father and of the Son and of the Holy Spirit in (\emph{ἐν}) flowing water (\emph{ὕδατι ζῶντι}). Having no flowing water, use another kind. If cold water is impracticable, take warm water. If you have neither, pour upon the head water three times in the name of the Father and of the Son and of the Holy Spirit.” The injunction is added: “Before baptism both the person that baptizes and he that is baptized shall fast, and such others as may be able to do so. The person baptized shall
\quellseite{46}
be commanded to fast for one or two days previous.” No more mention is made here of the laying on of hands than, for instance, in Acts 2:38 ff. The recitation of the Lord’s Prayer and the reception of the Holy Communion probably followed Baptism (ch. 8–10).

The Scriptures are silent on the question who imparted baptismal instruction. That the apostles and their fellow workers did most of this work is self-evident; not unlikely, however, they were assisted by laymen provided they were rich in knowledge. It is reasonable to suppose that the teachers mentioned beside apostles and prophets in 1 Cor. 12:28 and Eph. 4:11 f. did not only teach members of the congregation (Gal. 6:6), but also candidates for baptism. Paul taught large numbers at Antioch for a year (Acts 11:26), he taught at Corinth for 18 months (Acts 18:11), and he disputed at Ephesus for two years and three months (Acts 19:8 ff.); undoubtedly he instructed most of those who applied for baptism, no matter how great the number of his co-laborers. On the other hand, Aquila and Priscilla “expounded the way of the Lord more accurately to Apollos” (Acts 18:26) and evidently did this before his baptism. Certainly they were not the only “laymen” qualified for teaching; but gradually the “elders” and the above mentioned “teachers” assumed the responsibility for preparing the baptismal candidates. Apparently, itinerant teachers are mentioned in \emph{Didache} 12 if we may judge from statements in ch. 13.

The writings of the Apostolic Fathers do not add anything to the information so far gained. The “Two Ways” are worked into the \emph{Epistle of Barnabas} (ch. 18–20). Eusebius testifies (\emph{Hist. eccl.} iii 3, 6) that even in his day many still valued the \emph{Mandata} of Hermas as indispensable aids for instruction of catechumens. Justin does not offer anything essentially new; he writes (\emph{Apol.} i, 61): “All those who have
\quellseite{47}
been convinced and believe that what we teach and speak is true and promise to live accordingly, are induced to pray and to fast and to implore divine forgiveness of the sins of the past; and we pray and fast with them.” Thereupon baptism is administered.

\quellseite{48}
\section[§ 6. The Graded Catechumenate]{§ 6. The Development of the Graded Catechumenate between 180 and 325}

*G. v. Zezschwitz i, ch. 8–9; ii, 1, ch. 17. 21. 26. 27. 30–32. *E. Sachsse, pp. 4–115. E. Chr. Achelis ii, pp. 281–294. *J. H. Holtzmann, pp. 76–85. E. Cohrs, \emph{Katechumenat} in PRE³, 1901. O. Moe, \emph{Die Apostellehre u. d. Dekalog im Unterricht der alten Kirche}, 1896. E. v. d. Goltz, \emph{Das Gebet i. d. aeltesten Christenheit}, 1901. Concerning grades in the catechumenate see also: Fr. Probst, \emph{Lehre vom Gebet}, pp. 119 ff.; \emph{Katechese und Predigt}, pp. 39 ff.; \emph{Geschichte}, pp. 6 ff. F. H. Funk, \emph{Die Katechumenatsklassen des christlichen Altertums} (Theol. Quartalschrift) 1883, pp. 41 ff.; 1886, pp. 355 ff.; 1889, pp. 434 ff. F. H. Funk, \emph{Kirchengeschichtliche Abhandlungen}, 1897, pp. 209–241. K. Holl, \emph{Enthusiasmus u. Bussgewalt im griechischen Moenchtum}, 1898. Ed. Schwartz, \emph{Bussstufen und Katechumenatsklassen} (Schriften der wissenschaftlichen Gesellschaft Strassburgs vii), 1911. \emph{Irenaei opera}, ed. by Stieren, 2 vols., 1848–1853. \emph{The Five Books Against Heresies}, German by E. Klebba, KB², 1912; English in ANF, vol. ii; \emph{Presentation of the Apostolic Message} (\emph{Ἔνδειξις τοῦ ἀποστολικοῦ κηρύγματος}) German by Karapet Ter-Mekertschian and Edward Ter-Minassiantz, with notes by A. Harnack, ²1908; German also by S. Weber, Irenaeus ii, KB², 1912. \emph{Clementis Alexandr. opera}, ed. by O. Staehlin, 4 vols., 1905 ff. (BGSchr); Selected Writings, German by L. Hopfenmueller and J. Wimmer, KB¹, 1875; \emph{Protreptikos}, pp. 75–214; \emph{Paidagogos}, pp. 215–465. \emph{Protreptikos}, \emph{Paidagogos}, and \emph{Stromata} in English in ANF, vol. ii. Bratke, \emph{Die Stellung des Clemens Alex. zum antiken Mysterienwesen} (Theol. Studien und Kritiken), 1887. \emph{Tertulliani opera}, ed. by F. Oehler, 1851–1853; German by H. Kellner, 1882; English in ANF, vol. iii and iv; \emph{Tertullians katechetische Schriften}, by H. Kellner, KB², 1912. Catechetical discourses (perhaps enlarged) we have at least in \emph{De baptismo}, \emph{De spectaculis}, and in \emph{De poenitentia} (\emph{De bapt.} 20; \emph{De spect.} 1; \emph{De poenit.} 6; \emph{De idol.} 24). \emph{De poenitentia} and \emph{De pudicitia}, ed. by E. Preuschen, ²1910; \emph{De praescript. haereticorum}, ed. by E. Preuschen, ²1910; \emph{Adversus Parxean}, ed. by F. Kroymann (Krueger's Collection), 1907. \emph{Cypriani opera}, rec. W. Hartel, 1868–1871 (CSEL). \emph{The Three Books of Testimonia} are found in
\quellseite{49}
vol. i, pp. 35–184; \emph{De dominica oratione}, vol. i, pp. 267–294. Selected Writings, German by U. Uhl, J. Niglutsch, A. Egger, KB¹, 1879, 1892. \emph{The Three Books of Testimonies}, English in ANF, vol. v, pp. 507 ff.; \emph{Treatise on the Lord’s Prayer}, pp. 447 ff. Concerning the authenticity and purpose of the \emph{Testimonia} compare Glaue’s article in \emph{Zeitschrift fuer neutestamentliche Wissenschaft}, 1907. \emph{Origenis opera}, ed. by P. Koetschau, E. Preuschen, and E. Klostermann, 1889 ff. (BGSChr). \emph{Books Against Celsus}, vol. i, pp. 52–324; vol. ii, pp. 1–293; \emph{On Prayer}, vol. ii, pp. 294–403; \emph{De principiis}, vol. v, pp. 1–364. Selected writings, German by J. Roehm and J. Kohlhofer, KB¹, 1874–1876; \emph{On Prayer}, vol. i, pp. 21 ff.; \emph{Against Celsus}, vol. ii and iii. \emph{De principiis} and \emph{Against Celsus}, English by F. Crombie, ANF, vol. iv, pp. 239 ff.; iv, pp. 395 ff. \emph{Canones Hippolyti}, ed. by H. Achelis (Texte und Untersuchungen vi, 4), 1891. \emph{Didascalia}, translated into German from the Syriac by H. Achelis and J. Fleming, 1904. *\emph{Didascalia et Constitutiones Apostolorum}, ed. by Fr. X. Funk, 2 vols., 1905. \emph{The So-called Apostolic Constitutions and Canons}, German by F. Boxler, KB¹, 1874; English in ANF, vol. vii, pp. 348 ff. P. Drews, \emph{Der literarische Character der neuentdeckten Schrift des Irenaeus} (Zeitschrift fuer neutestamentliche Wissenschaft), 1907. *A. Harnack, \emph{Ueber den privaten Gebrauch der heiligen Schriften in der alten Kirche}, 1912; translated into English by J. R. Wilkinson, Crown Theological Library, vol. xxxvi, 1912.

The end of the second century marked a significant change in the life of the Church. Now the period of rapid general growth, but also of universal persecutions began. The first of these persecutions (under Septimius Severus, 202) was directed against those who had resolved to enter the Church, the catechumens. These conditions placed upon the Church if she would be true to herself, the duty of exercising more care than ever before in the reception of new members. The Church recognized this task and adjusted herself to the new status; she rejected such as did not seem able to meet the new demands, subjected all others to a lengthy period of probation and preparation, and took loving care of them. More stringent rules of admission into membership were adopted, and a graded catechumenate was developed.

\quellseite{50}
The recently (1907) discovered treatise \emph{Ἔνδειξις τοῦ ἀποστολικοῦ κηρύγματος} which was written by \emph{Irenaeus} after 190, clearly shows what painstaking methods were used in the preparation of prospective church members, at least at Lyons in Gaul. The course of instruction was based upon the tripartite confession of faith, it also offered a \emph{history} of divine revelation and of the economy of grace from creation to the entrance of the Israelites into Canaan adding brief references to Solomon and the prophets, and closed with a detailed presentation of Christ’s incarnation and redemptive work as a fulfilment of prophecy. Throughout the course, especially at the beginning and at the end, much stress was laid upon Christian \emph{morality}.

More information may be gathered from \emph{Tertullian’s} writings. In Africa new members were received in accordance with definite rules (cf. also the \emph{Canones Hippolyti}, a constitution which Hippolytus gave to his independent congregation at Rome about 220 A. D.); the catechumens formed a distinct order, the church watching over them with loving care; they were taught by means of sermons or lectures, and received a thorough moral training. Cyprian’s statements do not reveal anything additional except that his \emph{Testimonia} give eloquent proof of his sincere efforts at introducing the catechumens into the \emph{Scriptures}; this booklet is a collection of Bible passages gathered and arranged by him for a catechumen named Quirinus with the view of acquainting him with the Christian faith and the Scriptures.

Distinct grades of catechumens, however, are found in Alexandria. The passage in which \emph{Origen} discusses this \emph{matter} (\emph{Contra Celsum} iii 51), has indeed been variously interpreted, but the following conclusions seem warranted: (1) In Alexandria entrance into the catechumenate was preceded by an exploration and by a course of instruction. During this time the catechumen received general information concerning
\quellseite{51}
Christianity and he was more or less closely observed as to his sincerity and uprightness. We have no means of knowing whether this task was performed by officers of the congregation or by experienced laymen. (2) When the catechumen had satisfactorily passed through this preliminary stage, he was advanced into the catechumenate proper and began to attend the congregational services. But there were two orders of catechumens, the one embracing the newly admitted, the other those who had demonstrated by their life their sincere devotion to Christian ideals. The last stage naturally was of short duration; why should those already sufficiently tested be further delayed? Transition from the lower to the higher class was signalized by the symbol of purification (i.e., renunciation and exorcism), and from the higher into full membership, of course, by baptism. In this passage Origen does not give particulars concerning the material or the method of instruction, but stress is once more laid upon moral discipline. In several other passages he mentions “Moses and his law” as the means of moral instruction, however without having in mind specifically the Decalogue. He attaches profound importance to \emph{Bible reading} and strongly commends it. He advises catechumens to begin with the Old Testament Apocrypha which present less difficulties to the reader and abound in moral maxims; these books, he says, offer the milk required by immature believers. According to several passages (\emph{Contra Celsum} vi 10; \emph{Hom. in Iudices} v; \emph{De Princ. praef.}) the catechumens were given instruction concerning the baptismal confession and the Sacraments toward the end of the last stage.

In the order of baptism found in the \emph{Apostolic Constitutions} (vii 39 ff.) we would have a valuable witness for catechetical instruction as practiced at \emph{Antioch} if it really owes its origin to Lucian Martyr (311/2). In that case, the material used at Antioch was identical with that used at
\quellseite{52}
Lyons. It is especially noteworthy that also here Biblical History was employed for moral training. Apparently also the graded catechumenate was known at Antioch; the transition from the first stage to the second was signalized by prayer and laying on of hands.

At the beginning of the fourth century the graded catechumenate was a universally recognized institution. The preliminary stage attested by Origen has been abandoned; all applicants, after passing an examination, entered into the order of catechumens; after remaining here for some time they applied for baptismal instruction which was brought to a close, usually at Easter, with the Sacrament. How much time the catechumen spent in this course depended, of course, on the circumstances; toward the end of our period, two or three years were generally required.

\quellseite{53}
\section[§ 7. The Catechumenate in Perfected Form]{§ 7. The Catechumenate in Perfected Form, about 325–450}

*G. v. Zezschwitz i, ch. 8–10; ii, 1, ch. 17. 22–24. 27. 30–33. H. J. Holtzmann, pp. 66 ff. *F. Wiegand, pp. 1–176. E. Chr. Achelis, pp. 286–294. Compare also the bibliography given in ch. 4 and 6. K. Krawutzky, \emph{Die Katechesen des Augustin und die Vaterunserkatechesen des Augustin} (Neues St. Hedwigsblatt Breslau), 1875–77. *K. Knoke, \emph{Zur Methodik d. bibl. Geschichte} i, ²1878. O. Moe, \emph{Die Apostellehre und der Dekalog im Unterricht der alten Kirche}, 1896. K. Weiss, \emph{Die Erziehungslehre der Kappadozier}, 1903. W. Kroll, \emph{Alte Taufgebraeuche} (Archiv fuer Religionswissenschaft viii, suppl. pp. 27 ff.), 1905. E. Chr. Achelis, \emph{Der Dekalog als katechetisches Lehrstueck}, 1905. *P. Rentschka, \emph{Die Dekalogkatechesen d. Augustin}, 1905. F. J. Doelger, \emph{Der Exorzismus im altchristl. Taufritual} (Studien z. Geschichte u. Kultur d. Altertums iii, 1 and 2), 1909. J. Kunze, \emph{Die Uebergabe der Evangelien beim Taufunterricht}, 1909. G. Rietschel, \emph{Lehrbuch der Liturgik} ii, 1909. R. Reitzenstein, \emph{Von den dreierlei Fruechten d. christl. Lebens, eine fruehchristliche lat. Predigt} (Zeitschrift fuer ntl. Wissenschaft), 1914; of importance in this connection on account of its emphasis of the Decalogue; compare E. Seeberg, \emph{Eine neuaufgefundene lat. Predigt aus dem 3. Jahrhundert} (Neue Kirchliche Zeitschrift, 6 and 7, esp. pp. 506 ff.), 1914. MYSTERY CULTS: Chr. A. Lobeck, \emph{Aglaophamus de theologiae Graecorum causis}, 2 vols., 1829. G. v. Zezschwitz i, ch. 10. *N. Bonwetsch, \emph{Wesen, Entstehung und Fortgang der Arkandisziplin} (Zeitschrift fuer hist. Theologie, pp. 201–299), 1873. E. Rohde, \emph{Psyche. Seelenkult und Unsterblichkeitsglaube der Griechen}, 1890 (the Eleusinian mysteries, pp. 257 ff.; the Orphic teachings, pp. 395–428). *G. Anrich, \emph{Das antike Mysterienwesen in seinem Einfluss auf das Christentum}, 1894. N. Bonwetsch, \emph{Arkandisziplin} in PRE³, 1897. A. Mommsen, \emph{Die Feste der Stadt Athen im Altertum}, 1898 (pp. 204 ff.; 405 ff.). P. D. Chantepie de la Saussaye, \emph{Lehrbuch der Religionsgeschichte} ii, ³1905 (pp. 362 ff.). L. R. Farnell, \emph{The Cults of the Greek States} iii, 1907 (pp. 164–199). A. Dietrich, \emph{Eine Mithrasliturgie}, ²1910. *F. Cumont, \emph{Die orientalischen Religionen im roemischen Heidentum}, 1910. *R. Reitzenstein, \emph{Die hellenistischen Mysterien}, 1910. A. Jakoby, \emph{Die antiken Mysterienreligionen}
\quellseite{54}
\emph{u. d. Christentum}, 1910. *F. Cumont, \emph{Die Mysterien des Mithra}, ²1911. H. C. Sheldon, \emph{The Mystery Religions and the New Testament}, 1918. *M. Manitius, \emph{Geschichte der christlich-lateinischen Poesie bis zur Mitte des 8. Jahrhunderts}, 1891. \emph{Commodiani carmina}, ed. B. Dombart, CSEL, 1887; \emph{Instructiones}, pp. 3–112; \emph{Carmen apologeticum}, pp. 115–188; The \emph{Instructiones} in English by R. E. Wallis, in ANF, vol. iv, pp. 199–218. \emph{Iuvenci evangeliorum libri quattuor}, ed. I. Huemer, CSEL, 1891. \emph{Sulpicii opera}, rec. C. Halm, CSEL, 1866 (\emph{Chronica}, pp. 3–105). \emph{Sedulii opera}, rec. I. Huemer, CSEL, 1885 (\emph{Carmen paschale}, pp. 14–146; \emph{Opus paschale}, pp. 175–303). *\emph{Cyrilli Hieros. opera}, ed. in Greek and Latin by W. K. Reischl, 2 vols., 1848. 1860. (\emph{Catechetical Discourses} i. pp. 1–320, and ii, pp. 1–398); \emph{Catechetical Discourses}, German by J. Nirschl, KB¹, 1871; English by Gifford, NPNF², vol. vii. \emph{Gregory Nazianzen}, MSG, xxxv-xxxviii; Selected Writings, German by J. Roehm, KB¹, vol. i, 1874 (\emph{Discourse on Baptism}, pp. 48–91). \emph{Gregory of Nyssa}, MSG, xliv-xlvi; Selected Writings, German by H. Hayd and J. Fisch, 2 vols., KB¹, 1874, 1880 (The \emph{Λόγος κατηχητικός}, vol. i, pp. 123–205; \emph{Discourses on Prayer}, vol. ii, pp. 10–81; \emph{Discourse on Epiphany}, ib., pp. 276–299); \emph{Oratio catech.} and \emph{Orationes de precatione}, ed. with Latin translation by J. G. Krabinger, 1838, 1840. \emph{Basil the Great}, MSG xxix-xxxii; Selected Writings, German by V. Groene, vol. i, 1875 (\emph{Admonition on Baptism}, pp. 402–417). \emph{Chrysostom}, MSG xlvii-lxiv; Selected Writings, German by M. Schmitz, KB¹, 1879 (\emph{Two Discourses for Catechumens}, vol. iii, pp. 90–131); English in NPNF¹, vol. ix. \emph{Ambrose}, MSL xiv-xvii; Selected Writings, German by F. H. Schulte, KB¹, 1871 (\emph{De mysteriis}, pp. 119–225). Ambrose is indirect author of \emph{De sacramentis}. The \emph{Explanatio Symboli} reprinted in P. Caspari, \emph{Ungedruckte, unbeachtete und wenig beachtete Quellen z. Geschichte d. Taufsymbols u. d. Glaubensregel} ii, 1869, pp. 48–127, and in P. Caspari, \emph{Alte und neue Quellen}, etc., 1879, pp. 196–222, is Ambrose’s own. The \emph{Exhortatio sancti Ambrosii episcopi ad neophytos de symbolo}, reprinted in Caspari, \emph{Ungedr., etc. Quellen} ii, 128–182, and in Caspari, \emph{Alte und neue Quellen} pp. 186–195, was not written by him; perhaps it was written by Eusebius of Verrelli and was delivered between 340 and 360. *\emph{Rufinus}, MSL xxi; \emph{Commentarius in symbolum apostolorum}, with notes and English translation, 1908; German by H. Bruell, KB¹, 1876 (pp. 19–85). \emph{Silvia vel potius Aetheriae Peregrinatio ad loca sancta}. CSEL xxxix, ed. W. Heraeus, 1908. *\emph{Augustine}, MSL xxxii-xlvii; \emph{De catechizandis rudibus}, ed. by G. Krueger, ²1909; in German by Th. Ficker, with introduction
\quellseite{55}
by G. v. Zezschwitz, 1863; in English by Shedd, NPNF¹, vol. iii. \emph{Enchiridion}, ed. by O. Scheel, 1903; German by J. Molzberger, KB¹, vol. iv of Augustine’s writings, pp. 572–695; English in NPNF¹, vol. iii. \emph{De fide et symbolo} and \emph{De fide et operibus}, CSEL, ed. by J. Zycha, 1900, pp. 1–97; English in NPNF¹, vol. iii. \emph{De fide et operibus}, German by R. Storf, KB¹, vol. iv of Augustine’s writings, pp. 485–565. \emph{De symbolo ad catechumenos}, German by R. Storf, ibid. pp. 351–481. Of this series of sermons only the first one was written by Augustine (2–4 were written by another African bishop, probably Vigilius of Tapsus); it is a sample of Augustine’s explanation of the Symbol before competentes. *In \emph{Sermo} 220, 228, 227 (MSL xlvii) we have samples of addresses delivered in connection with the \emph{traditio symboli}, or after baptism before the assembled congregation. The last mentioned sermons are reprinted in Lietzmann, \emph{Ausgew. Pred.} ii (Kleine Texte 13), 1905, pp. 10–16. Cf. Wiegand, pp. 42 ff.; 83 ff. \emph{Petrus Chrysologus}, MSL lii; Selected Discourses, German by M. Held, KB¹, 1874 (esp. *\emph{Three Discourses on the Symbol}, pp. 58 ff.; 313 ff.; 321 ff.; and *\emph{Three Discourses on the Lord’s Prayer}, pp. 65ff.; 326ff.; 331ff.); cf. Wiegand, pp. 120ff. \emph{Niceta of Remesiana}, MSL lii; \emph{Explanation of the Symbol} in Caspari, \emph{Kirchenhistorische Anekdota} i, 1883, pp. 341–360; cf. Wiegand pp. 108ff., and A. E. Burn, \emph{Niceta of Remisiana, His Life and Works}, 1905. \emph{Maximus of Turin}, MSL lvii; \emph{Tractatus de baptismo}, an explanation of the baptismal rite for Neophytes, pp. 771–782; there we find also an address on the symbol delivered in connection with the “tradition” of the same; cf. Wiegand, pp. 133ff.

The graded catechumenate was developed by the Church during a period of bitter persecution and severe adversity. By 325 it was universally recognized and established. And although now oppression and tyranny ceased, the Church possessed sufficient vitality to keep this institution efficient for over a hundred years, surely a commendable achievement. The preliminary stage mentioned by Origen had indeed been reduced to a mere introductory catechesis, but everywhere throughout the Church a distinction was now made between catechumens and \emph{φωτιζόμενοι} or applicants for baptism.

Augustine’s work \emph{De catechizandis rudibus} supplies the best information concerning the mode of reception into the
\quellseite{56}
\emph{catechumenate}, and the introductory catechesis. Upon the request of Deogratias the deacon, he does not only theoretically discuss the subject matter and the method of the introductory catechesis, but he also offers two model catecheses, one long, the other short. According to this treatise, a heathen who at that time (about 400) intended to join the Church, was examined as to his motives by the bishop or a presbyter or a deacon. In all probability inquiries in regard to the applicant were made also among others, for the number of nominal Christians who had joined for social and business reasons was constantly increasing. If ignoble motives were discovered, an attempt was made to treat the applicant individually and so to implant purer ones in him. The catechetical lecture—it did not proceed in the form of questions and answers; cf § 1—was, of course, adapted to the special needs of the individual applicant; but in all cases it was composed of two elements: a \emph{historical survey of the kingdom of God} from creation to the present, and the inculcation of the \emph{mandates of true morality}. Employing the argument from prophecy and fulfilment, the catechist would present Jesus and his Church as the center of history; he would point out that the infinite love of God is the all-impelling motive and that, in return, God expects the Christian to love him and his fellow-man; and he would endeavor to re-enforce his admonitions to true morality by turning the applicant’s mind to the last judgment and life eternal.

The introductory catechisation finished, the applicant avowed his faith in what he had heard, and his intention to pattern his life after the commandments received. Several liturgical acts were performed by which he was “made” a \emph{catechumen} or Christian, and thenceforth he was regarded as a member of the congregation. For the time being he received no further formal instruction. By conscientious attendance of public worship (\emph{missa catechumenorum}; Scripture reading
\quellseite{57}
and sermon) and by fruitful association with experienced Christians, it was thought, he would be gradually introduced more thoroughly into Christian faith and life. Writings like those of Commodian (\emph{Instructiones} and esp. \emph{Carmen apologeticum}; middle of third century), of Sulpicius Severus (\emph{Chronica}; about 400), of Iuvencus (\emph{Libri evangeliorum}; about 330), and of Sedulius (\emph{Carmen paschale}; about 400–450) were well suited to acquaint the catechumens more thoroughly with sacred history and the progress of the Church. In the West, the catechumens were often called “hearers” because of their obligation to attend the church services. From \emph{Apost. Const.} viii, 6 it appears that by a special petition the congregation remembered its catechumens every Sunday. Throughout this period, instruction in true morality was kept in the foreground.

Unfortunately an injurious practice arose during this period: many of the catechumens remained for years and decades in the first grade of the catechumenate and did not enter the second grade until they were old and gray, frequently not until they were on their deathbed. Various reasons may be adduced for this aberration. Ever since 325 the number of applicants had so tremendously increased that the Church frequently was unable to exercise the necessary discrimination at their reception into the catechumenate. Still more catechumens, however, voluntarily remained in the lower class because it was taught that certain sins done after baptism were unpardonable or that forgiveness for them might be attained only upon the arduous way of penitential discipline. Thus many sincere souls postponed their baptism in order not to jeopardize their salvation. But much more frequently this was done by frivolous and light-minded persons who had entered the catechumenate for the most superficial reasons and continued to live in sin; they refused to submit to church discipline, yet desired to belong to the church and eventually
\quellseite{58}
to partake of salvation. Baptism administered shortly before death appeared to them as a safe and convenient indulgence for a life of sin. This sinful practice continued to spread in spite of the earnest efforts of many faithful preachers. In this way, the first grade of the catechumenate became so protracted that many catechumens were totally unknown to the bishop of the congregation when they entered into the second grade; real baptismal instruction under these circumstances did not begin until then. Nevertheless, the idea of successive stages was still retained for some time.

The catechumen who desired to be baptized, had to announce his intention and his name to the bishop at the beginning of lent. If his previous life presented no obstacles, he was duly matriculated. This moment marked the beginning of a very significant period of his life. As \emph{baptismal candidate} (\emph{φωτιζόμενος}, competens) he stood in much closer relation to the congregation which, in turn, devoted itself to him with particular interest, directed his thoughts to the great day of Baptism, and endeavored to prepare him inwardly for it.

Three sets of influences were brought to bear upon the candidate: \emph{ascetic, educational, and liturgical}.

Already in the preceding periods candidates for baptism had been \emph{ascetically influenced} — we use this word in the widest sense. Now, however, this was done to a much larger extent, possibly in reverse proportion to the candidate’s personal preparation. Confession of sin, required ever since apostolic times, was retained, and the “competent” period was represented as the time best suited for confessing although an enumeration of all sins was not demanded. Fasting formerly practiced during the time immediately preceding baptism, was extended over the whole period. The competentes were expected to show their sincerity in many ways: they must not attend theatres or banquets, married people must live
\quellseite{59}
continentally, all of them must offer daily prayer with genuflection, read the Scriptures daily, and give alms.

The bishop himself or, by his order, the presbyter imparted the prebaptismal instruction. As far as the East is concerned, we receive the fullest information about this through the \emph{Catecheses} of Cyrillus of Jerusalem, i.\,e., discourses delivered by him to \emph{competentes} in 348 when he was presbyter at Jerusalem. These lectures may well be considered “the typical catechetical sermons of antiquity.” An exhaustive explanation of the apostolic symbol forms the center of gravity. 13 catecheses out of 18 are devoted to it. The third lecture discusses baptism. But the ethical element was not absent; it appears now and then in the more dogmatical and polemical sections on the symbol, and with special emphasis in the preliminary and the first two catecheses. Chrysostom pursues the same aims in the two extant discourses for \emph{photizomens}. He wishes to inform them thoroughly of the nature and effect of baptism and to arouse them to a life in harmony with the Christian confession.

The best witness for the instruction of \emph{competentes} in the West is Augustine whose testimony is corroborated and supplemented by a number of others. No little importance was attached to moral training. In \emph{Sermo 216} Augustine, as young presbyter (391), charges the \emph{competentes} whom he was addressing: “Do not withdraw your hand from the plow! Away from the world! Be God’s own!” A compendium of moral instruction did not exist. Not even the Decalogue was used as such, although Augustine presupposed its knowledge and occasionally referred to it. The practical needs of the \emph{competentes} suggested the ethical subjects to be discussed, and texts were chosen accordingly if the lectures were based on texts at all. All precepts, even the 10 Commandments, were generally summed up in the “Twin command of love” as the sum and substance of all true morality. The trend of the
\quellseite{60}
times did not favor such moral training; what was demanded was an easier approach to baptism, and restriction, yes, even elimination of moral training before baptism. The prevalent moral laxity and the practical impossibility of exercising adequate moral control over thousands of catechumens served to obscure the Church’s sense of responsibility; and there were many leaders who argued that \emph{photizomens} should receive doctrinal instruction only and that moral instruction should not be imparted until after baptism. Even adulterers were admitted to the second grade and to baptism, and this practice was defended in all seriousness; it was thought that only idolaters and men of disgraceful occupation could be excluded from the Sacrament. Augustine’s forceful book \emph{De fide et operibus} (413) is directed against these demoralizing tendencies, and he himself clung to the moral training of \emph{competentes}; he succeeded in retarding the movement, but was unable to conquer it.

Also Augustine accorded to instruction in the symbol more space than to the other subjects. This is demonstrated by his short sermons upon the symbol (\emph{Sermo 213–215}) and esp. by the tract \emph{De symbolo ad catechumenos} (MSL xl, pp. 627–636). The former were delivered when the \emph{competentes} received the symbol or recited it (see page 42), and the latter is an exhaustive explanation of the symbol designed as a means of further instruction of \emph{competentes}. Others accord similar eminence to the symbol. Rufinus, the erudite presbyter of Aquileia, for instance, published in his \emph{Expositio in symbolum apostolorum} a suitable manual for the use of bishops and presbyters in the instruction of \emph{competentes}. Again, Peter Chrysologus, the indefatigable bishop of the rising city of Ravenna (430–451), explained the symbol to \emph{competentes} as is evident from his sermons. Maximus of Turin testifies to the prevalence of the same usage in northern Italy. Likewise Nicetas, the missionary bishop of Dacia
\quellseite{61}
(about 400–410), explained the symbol to his \emph{competentes} as may be seen from an extant address which he delivered when they recited the symbol. It is not unlikely that the three catecheses on the symbol generally found in connection with Augustine’s \emph{De symb. ad. catech.} (\emph{Aug. Serm. 237–239}; MSL xxxix, 2183 ff.) are of African origin (possibly Vigilius of Tapsus is the author). This impressive array of names permits of only one conclusion: Whatever stress was laid upon ethical instruction, summarized in the twin command of love (cf. in addition to Augustine esp. Nicetas), the symbol was the article that overshadowed all others; it was the real verbum abbreviatum. The early Church thought that every subject of faith and doctrine was summarized in the symbol.

At some places in the West also the \emph{Lord’s Prayer} was an integral part of the instruction imparted to the \emph{competentes}. In Africa, for instance, Augustine delivered it to them “within the \emph{quadragesima},” i.\,e., on the Saturday before \emph{Judica} so that they might commit it to memory and at their first communion pray it with the congregation. He also gave them a concise explanation so that we find brief catecheses on Prayer corresponding to those on the symbol. Chrysologus of Ravenna followed the same custom, as his discourses on the Lord’s Prayer show.

The most important characteristic of this period, however, was a feature adapted from paganism, \emph{arcane discipline}, with which a multitude of \emph{liturgical acts} was connected. It is well known that in the second and third centuries the traditional cults of Greece and Rome rapidly decayed and in ever increasing manner failed to supply men’s religious needs. As a result, some of the ancient Graeco-Roman mystery religions, particularly the Eleusinian cult, were revived, and a number of Eastern mystery cults, notably Orphism and Mithraism, became widely diffused in the mediterranean world. Surprisingly large numbers entered the societies of devotees
\quellseite{62}
in their quest of redemption and immortality. The Christians were in daily contact with these mystery worshippers; they heard them boast that the mystery cult with its washings, processions, symbolic acts, with its secret formulae and rites, gradually leads to the vision of deity and atonement and that it secures the hope of immortality. It is only natural that converts from the ranks of cultured paganism were tempted to reply in defense: “We possess the true mysteries, Baptism and Communion. Fellowship with God is attained only through them,” and that they made efforts to surround the Sacraments with a halo of secrecy and to embellish these simple acts with a mass of liturgical acts. The Church unfortunately did not long resist these temptations, but appropriated the fundamental features of the mystery cults and thus introduced the so-called arcane discipline. It must be remembered, however, that by acceding to these demands the Church not only met the trend of the times, but also the hierarchic imperialism which since the end of the third century rapidly became conspicuous. Now the priest became the mystic-theurgic agent on whose performances the validity of all sacred rites depended, the priest’s activity was indispensable, without him there was no approach to God. At the beginning of our period this transformation had been accomplished. Baptism and communion now are known as \emph{μυστήρια}, baptismal preparation as \emph{μύησις}, the officiating bishop as \emph{μυσταγωγός}, and the baptized as \emph{μεμνημένοι} or \emph{ἐπόπται}. Just as in the Eleusinian mysteries the climactic feature of the initiation, the opening of the temple, took place at night, so also Baptism was performed at night, preferably in that of the Passover. Just as a series of ceremonies prepared the candidate for the initiation into the mystery and for the vision of deity, so now various holy acts preceded baptism. Strict secrecy must necessarily be observed in regard to the “mysteries” of Baptism and Communion; a little later, Christians
\quellseite{63}
were pledged to secrecy even regarding the preparatory rites and particularly the symbol and the Lord’s Prayer. This is the explanation of the fact that the symbol itself was delivered in its integrity in a final solemn act, although the \emph{competentes} had received it piecemeal during the period of instruction, and that they were not given the Lord’s Prayer until immediately before baptism.

Thus the following series of liturgical acts resulted: Having heard the introductory catechesis, the applicant in general terms expressed his faith and his intention to obey the commands of the Church. Now he was “made a Christian” by the \emph{signatio crucis} upon the forehead, the imposition of hands, and (in the West) the exorcised and consecrated salt (\emph{Christianum facere, catechumenum facere}). Especially the \emph{datio salis}, was in the West, accounted as \emph{sacramentum catechumenorum}. When he entered into the second grade, his name was entered in the roster of the congregation by the bishop or a presbyter, and in some instances he was given a new name (\emph{nomen dare}). This was followed by exorcism, i. e., the evil spirit was expelled by the laying on of hands, afflation, and adjuration, a procedure based on the notion that pagans, like the possessed, were physically, no less than mentally and spiritually, in the power of Satan. By this act of exorcism the candidate was enrolled as a competent; the second grade of the catechumenate was supposed to be the period during which the soul is delivered from Satan’s power and dominion. In the West, another act, the \emph{apertio aurium} (cf. Mark 7:33), was performed at this moment; the bishop touched the candidate’s ears and nose (not the mouth; occasionally exorcised oil was used in this rite) and said: “Ephphatha!” to indicate that henceforth he would be able to understand the words of the gospel and to perceive “the sweet savor of Christ.” In some localities this symbolism was carried out still further by the \emph{traditio}
\quellseite{64}
or expositio evangeliorum: after a short discourse on the term “gospel”, on the four evangelists, and on the symbols generally assigned to the Gospels (\emph{homo, leo, vitulus, aquila}, cf. Eze. 1:10), the bishop read the first verses of each of the four Gospels. The candidate’s response to these acts was a solemn renunciation. From now on he was aided in his battle against the forces of evil by a number of “scrutinies”, i.\,e., religious meetings at which various liturgical acts were performed on him, as the sign of the cross, imposition of hands, exorcisms, and other supposedly beneficial ceremonies. The climax was reached with the \emph{traditio symboli}, when the symbol, this secret, wonder-working formula, was delivered to him in a solemn church service; as a rule, a sermon on the symbol was preached at this occasion (cf. August. Serm. 212–214). In Africa this was done on the Saturday before Laetare, at other places usually on Palm Sunday. The next step corresponded to the traditio; it was the \emph{redditio symboli}, i.\,e., the recitation of the symbol by the competent. As a rule, this was done one week after he had received it (Saturday before Judica, Saturday before Easter, Maundy Thursday, in Jerusalem on Palm Sunday); profound importance was attached to this act of confession before the assembled congregation. No less significant was the \emph{traditio orationis dominicae}; the \emph{competentes} were expected to commit the Lord’s Prayer to memory so as to be able to pray it with the congregation before Communion. In Rome two scrutinies were conducted on Good Friday and the final scrutiny on the following day. Finally, the candidate was sufficiently well prepared for baptism.

This whole procedure impresses us as rather grotesque and strange. Discontented with the divinely instituted means of grace, men resorted to all sorts of human devices and crowded back indoctrination in the essentials of Christianity. So much is certain that the hierarchy welcomed this mummery;
\quellseite{65}
for the individual was now hopelessly chained to the church whose bishops and priests were the only persons capable of performing these acts that purified the soul and conveyed grace, that paved the way to baptism and thus to communion with God.

\quellseite{66}
\section[§ 8. Decay of the Catechumenate]{§ 8. The Decay of the Catechumenate (about 450–600)}

G. v. Zezschwitz i, ch. 9. J. H. Holtzmann, pp. 86 ff. *F. Wiegand, pp. 181 ff. F. Probst, \emph{Katechese und Predigt vom Anfang des 4. bis zum Ende des 6. Jahrhunderts}, 1884. Wilson, \emph{The Gelasian Sacramentary}, MSL lxxiv, 1894.

The catechumenate of the Early Church received its first blow when the heathen in large masses crowded to the baptismal font. Still, a considerable measure of Christian knowledge was imparted and a reasonably thorough moral training was accomplished as long as the Church insisted upon a careful examination of the baptismal candidates, upon a catechumenate of sufficient length, and, especially, upon thorough instruction during the competent period. Now, however, just these three essential conditions were increasingly disregarded. The preparatory catechetical discourse was discarded, probably as early as the fifth century. The time of the catechumenate was dangerously shortened. Thus in 506 the Synod of Agde declared that Jewish converts must remain in the catechumenate for 8 months; hence, ordinarily the time of preparation was still briefer (this action was taken because experience had showed that many former Jews withdrew from the Church soon after they had joined). This, however, was not all. At the council at Bracara in 610 the time for the instruction of competentes was actually cut down to 20 days; and in most cases religious instruction was so completely overshadowed by ever increasing scrutinies that only a few formulae were memorized. That was all that remained of instruction which at one time had been conscientiously cultivated and cherished as an indispensable obligation of the Church. But all the alien elements adapted from paganism,
\quellseite{67}
the ceremonies and magic formulae, were retained in the Church. No wonder that this period produced large numbers of sacramentaries and liturgical treatises and books about the scrutinies, but very few catechetical writings. The catechumenate decayed, it was buried under the scrutinies. This would perhaps have been less dangerous if conditions had been such that but few unbaptized adults remained within the precincts of the Church; but the very opposite is true: as the Roman Empire crumbled before the onslaught of the barbaric peoples, ever new pagan nations settled within the boundaries of the Church which needed thorough instruction and training. The sad state of affairs was aggravated by the fact that even those baptized in infancy received no regular or sufficient training. A spiritually sterile church, poisoned by hierarchic thought with its emphasis on mystic-theurgic acts, was unable to renew these nations inwardly though she subdued them outwardly—with the help of the state.

\quellseite{68}
\section[§ 9. Conclusion of the Catechumenate]{§ 9. The Conclusion of the Catechumenate in the Early Church}

J. W. Hoefling, \emph{Das Sakrament der Taufe}, vol. i, 1859, pp. 277–550. G. v. Zezschwitz i, ch. 9. E. Sachsse, pp. 74–92. *F. Wiegand, pp. 217ff. *G. Rietschel, \emph{Lehrbuch der Liturgik} i, pp. 13–58. P. Drews, \emph{Taufe, Liturgischer Vollzug} in PRE³, 1907. Compare the literature for chapters 7 and 8.

The goal of instruction in the Early Church was baptism; this holy act formed its conclusion. As early as the second century (in Rome, e.\,g., in 194), Easter Sunday was thought to be the most appropriate time for baptism, in as much as the baptized “were buried with Christ into death and like as Christ raised up from the dead.” Tertullian, however, mentions Pentecost beside Easter and admits that in itself any day is suitable for baptism. The Synod of Carthage, 256, concurred in this view, at least in respect of the baptism of infants. In the fourth century the custom arose to baptize all competentes jointly at Easter and—save in imminent danger of death—not to baptize at all except at stated seasons. The Apostolic Constitutions, Basil, Chrysostom, for instance, prescribe the night between Easter eve and Easter morn for baptism. Thus the act became very impressive: the somber period of lent just ending, the stirring final appeals of the instructing bishop, the serious vows soon to be offered, all served to enhance the solemnity of the hour. However, alongside of Easter, Pentecost retained its baptismal celebration; at some places baptisms took place on the festival of Epiphany or even on Christmas Day (celebrated on December 25th as early as 354 in Rome, at Constantinople in 379, at Antioch in 388). When adults were no longer baptized, Easter was stripped of its chief solemnity
\quellseite{69}
except where the children born during the year were baptized on that day.

Formerly every place was deemed appropriate for baptism; Tertullian goes so far as to say that it makes no difference whether one is baptized in the ocean, in a pond, in a river, in a fountain, or in a basin. However since baptism was performed in the night, special baptisteries were needed. They were either added to the church structure or built in its vicinity.

According to mode, baptism was performed by immersion; but this was by no means now, nor had it ever been, the exclusive form. It is not likely that the 3000 converts added on Pentecost Day were immersed. The passage from the \emph{Didache} quoted above (cf. p. 23) explicitly permits another form of baptism not only for the sick, but in all cases where there was not sufficient water for immersing. Tertullian recognizes the validity of “baptism in a fountain.” In the cemetery of Callistus (217–222) in Rome two illustrations of the baptismal act are found, both representing the candidate for baptism standing naked in water which reaches to his ankles; according to the one the baptizer pours water over his head, according to the other over his head and body.

While in the earlier times baptism was a very simple act, an extensive ritual of baptism was developed later. As early as the fourth century the following ceremonial was generally observed: When the last scrutiny had been completed on the great sabbath, the baptismal act commenced at about three o’clock in the afternoon. In solemn procession the whole group marched to the church; here the candles were consecrated symbolizing the φωτισμός, and several prayers and scripture lessons were read, generally concluding with Psalm 42. Now the procession moved to the baptistery. In the meantime night had fallen and frequently the whole
\quellseite{70}
city was illuminated for the occasion. Again several liturgical acts were performed before baptism was administered. Cyril of Jerusalem has given us the most graphic account of these ceremonies: The candidates for baptism entered the vestibule of the baptismal chapel and took their position on the West side. Stretching forth their hands, they renounced \emph{Satan and all his works, his pomp and his dominion} (\emph{abrenuntiatio diaboli}, \emph{ἀποταγὴ τοῦ Σατανοῦ}). Now they shifted to the East side of the chapel and confessed, one by one: “I believe in the Father, and in the Son, and in the Holy Spirit, and in the one baptism of repentance” (\emph{συνταγὴ τοῦ Χριστοῦ}). In the West this confession was made in the form of questions \emph{and answers} and therefore called \emph{interrogatio de fide}. After a separation of the sexes the candidates proceeded into the baptistery proper, divested themselves of their garments, and were \emph{anointed} from head to foot with sacred oil (in the West only upon certain parts of the body). This anointing conveyed supposedly heavenly powers. Then the baptismal water was consecrated through invocation of God and the sign of the cross. And now finally the baptismal act was performed. The candidates descended into the water and were thrice immersed in the name of the triune God (in Spain but once, as a symbol of the unity of the three persons of the Trinity, and a testimony against the Arians). Occasionally the several parts of the Creed were recited, the candidate replying after each, “I believe,” before he was immersed. The baptizer concluded with a supplication like the following: “God the almighty Father, who regenerated thee of water and the Spirit, and forgave all thy sins, anoint thee unto life eternal.” Therewith forgiveness of sins had been bestowed upon the candidates and they were received as children of God.

Three additional acts concluded the ceremony: another anointing, the making of the sign of the cross, and the
\quellseite{71}
imposition of hands. Especially the last was highly esteemed particularly in the West and already in Jerome’s times it was performed only by the bishop and frequently in a special room, the \emph{consignatorium}. Soon these supplementary acts were valued even higher than baptism because the impartation of the Spirit (cf. Acts 8 and 19) and of His gifts (cf. Isa. 11:2) was ascribed to them (beginnings of the \emph{confirmatio}). The newly baptized were arrayed in white garments as symbol of purity and innocence; at times a mixture of milk and honey was given them as a concluding act to indicate that at last they had arrived in the Church, the true Canaan. The climax of all acts was the first Communion. As believers and full members of the Church they enjoyed on Easter Sunday for the first time the privilege of attending the \emph{missa fidelium}, of praying the Lord’s Prayer with the assembled children of God, and of approaching the Lord’s table.

A brief period of supplementary training followed upon baptism. As far as the East is concerned, this practice is attested by the mystagogic catecheses which Cyril of Jerusalem addressed to the neophytes in the week after Easter. He explained to them the various ceremonies connected with baptism (Catecheses 1–3; the nature and meaning of baptism itself had been discussed during the competent period), also the nature and blessing of the Lord’s Supper (Cat. 4), and of the liturgical acts embellishing the same (Cat. 5). In the last catechesis he also explained the Lord’s Prayer because this prayer spoken with a loud voice by the whole congregation immediately before approaching the altar formed part of the liturgical elaboration of the Eucharist, and because he never explained it before baptism. Also in the West supplementary instruction concerning the Sacraments was imparted in the week after baptism, and if the Lord’s Prayer had not been explained before, this was done now. Augustine’s Sermones clearly show that earnest efforts were made to awaken
\quellseite{72}
in the \emph{renovati} or \emph{infantes} the firm resolution to lead a new life. Frequently, but by no means always, the Decalogue was used in this connection.—During this week the newly baptized were accorded special attention also in the public services, and they were allowed to appear in their white baptismal vestments. They had to take them off again on the Saturday or Sunday after Easter which in consequence was known as \emph{Dominica post albas} or \emph{Dominica in albis} (the latter perhaps an abbreviation for \emph{in albis depositis}). The neophytes were expected not to attend banquets, etc., but fasting was not required: a season of gladness had begun.

\quellseite{73}
\section[§ 10. Infant Baptism and Instruction]{§ 10. Infant Baptism and the Christian Instruction of the Young in the Early Church}

J. W. Hoefling i, pp. 98–126. *G. v. Zezschwitz i, ch. 13 and 19. *E. Sachsse, pp. 92–101. *F. Wiegand, pp. 217–238. Moeller-Schubert, \emph{Lehrbuch der Kirchengeschichte} i, ²1902, pp. 806ff. G. Thomasius, \emph{Dogmengeschichte} i², pp. 547 ff. A. Harnack, \emph{Lehrbuch der Dogmengeschichte} i⁴, pp. 472ff. G. Rietschel, \emph{Liturgik} ii, pp. 6–8. P. Drews, \emph{Taufe. Kindertaufe. Pateninstitut} in PRE³, 1907. J. Steinbeck, \emph{Lehrbuch d. kirchl. Jugenderziehung}, 1914, pp. 8–9. S. S. Laurie, \emph{Pre-Christian Education}, 1900. G. Rauschen, \emph{Das griechisch-roemische Schulwesen}, 1901. J. Brunner, \emph{Der heilige Hieronymus und die Maedchenerziehung}, 1910. A. Harnack, \emph{Ueber den privaten Gebrauch der heiligen Schriften}, 1912. M. Reu, \emph{Der kirchl. Jugendunterricht in der Alten Kirche}, in Kirchl. Zeitschrift, 1913. \emph{Cypriani Epistolae}, CSEL³, ii; English in ANF, vol. v; German by J. Niglutsch and A. Egger, KB¹, 1879; \emph{Letter to Fidus}, pp. 328ff. \emph{Basilii opera}, MSG xxix-xxxii; \emph{Regulae} in vol. xxxi; \emph{Selected Writings} in English in NPNF², vol. vii; German by V. Groene, KB¹, 1877; \emph{55 regulae}, pp. 37–161, in the fifteenth of which we find four chapters on reception and education of children in convents (pp. 85ff.), cf. \emph{regula 53}. *\emph{Sancti Ioannis Chrysostomi de educandis liberis liber aureus}, ed. by F. Combesis, 1656; English: \emph{The Golden Book of St. John Chrysostome, Concerning the Education of Children}, translated by John Evelyn, 1658 (reprinted in: \emph{The Miscellaneous Writings of John Evelyn}, 1825); German: \emph{Ueber Hoffart und Kindererziehung samt einer Blumenlese ueber Jugenderziehung aus des Chrys. Schriften}, ed. by S. Haidacher, 1907.

The Catechumenate of the Early Church was an institution designed for the preparation of adults for membership in the Church. However, what was the status of infant baptism and of the instruction of children in the Early Church?

That infants stand in need of baptism, the New Testament clearly states when it teaches original sin and hereditary
\quellseite{74}
guilt (John 3:6; Rom. 5:12. 18. 19; Eph. 2:3); and when examined in the light of the context, 1 Cor. 7:14 confirms, rather than contradicts, the necessity of infant baptism. That infants may be baptized and that Christ would bless them, is attested by Luke 1:15, by Christ’s own course of action, by Christ’s direction to the disciples (Matt. 19:14), and by the universality of the divine redemptive will (1 Tim. 2:4; Matt. 28:19).

It is quite another question whether the Church universally recognized the duty of infant baptism and fulfilled its obligation. From Eph. 6:1 it would appear that grown up children were baptized, for the apostle addresses the children, no less than women, men, and slaves, as members of the congregation; evidently he presupposes that children no less than their parents had been baptized. No case of infant baptism is explicitly recorded in the New Testament, but events as those narrated in Acts 16:15. 33; 1 Cor. 1:16, and the significant type of circumcision very strongly suggest that it was practiced in apostolic times. The absence of explicit attestation undoubtedly is due to the fact that the Early Church, like Foreign Missions today, concentrated its energy upon the winning of adults, and that infant baptism could not possibly become a general rule until there were Christian homes which would guarantee subsequent training. Moreover, that infant baptism was actually, even if not universally, practiced, is borne out by the manner in which it is discussed by its first undisputed witnesses. They do not speak of it as a recently adopted custom, but as a familiar and widely prevalent usage. Compare, e.g., the words of Irenaeus in his \emph{Books Against Heresies} ii, 22 (written during the reign of Commodus, 180–192): “Christ came to save all through himself—all, I say, who are regenerated of God: sucklings and children and boys and young men and old people,” and bear in mind that according to Irenaeus’ usage regeneration
\quellseite{75}
refers to baptism (\emph{Adv. Haer.} iii 17, 1). In the same manner Tertullian testifies to the practice of infant baptism when he warns against a premature baptism of children which practice—to judge by his words—must have been justified by an appeal to Matt. 19:14 (\emph{De bapt.} 18). Still more plainly Origen actually asserts: “\emph{Ecclesia ab apostolis traditionem suscepit etiam parvulis baptisma dare}” (\emph{In Rom.} v 9), and defends infant baptism by referring to the natural pollution of infants (\emph{In Luc. evang. hom. 15}; \emph{In Levit. hom. 8}). At the synod of Carthage, 256, the question was not whether infants should be baptized at all, but whether they should be baptized before the eighth day. Cyprian’s letter in this matter to Fidus, the presiding bishop, exhibits surprise that any indecision concerning this question existed. In the \emph{Constit. Apost.} the injunction is found: “Moreover baptize also your children (νήπια) and bring them up in the nurture and admonition of the Lord!” (vi 15), and in viii 11. 13, also in Cyprian’s writings, the communion of children is mentioned which postulates their baptism. It is true that in the third century parents now and then had to be admonished to have their children baptized, and that in the fourth century even pious parents let their children grow up without baptism; but this does not at all mean that infant baptism had not yet met universal recognition, it is rather due to the fear that all the blessings received in baptism would be forfeited by subsequent sins, and undoubtedly also to the unwillingness of assuming the parental obligations inseparably connected with infant baptism (Cf. esp. Tertullian De bapt. 18 and Apost. Constit. vi 15). Gregory Nazianzen whom the parents did not have baptized in his youth, grows emphatic in his demand that children should be baptized immediately when danger threatens, otherwise in the third year (\emph{Discourse on Baptism} 17, 28). Augustine speaks of countless numbers of baptized infants, and Leo the Great, in an Easter sermon, of thousands of infants
\quellseite{76}
snatched out of Satan’s power. And Gregory the Great imposes triennial penance upon parents who have permitted a child three years old to die without baptism. Usually infants were baptized at the seasons fixed for the baptism of adults.

A grievous error was committed when the complete ritual of adult baptism was employed in the baptism of infants. Without abridgment the whole series of liturgical ceremonies was performed upon the children, the responses, especially the renunciation and the confession, being made by the sponsors. Such seems to have been the condition already at Tertullian’s times (\emph{De bapt.} 18). So it happened that the act of renunciation, originally presupposing idolatry, was insisted upon even in the case of children of Christian parents; not even its form was changed: “Do you renounce? Do you believe?”, although Augustine already said that these baptismal questions could no longer be understood in their proper sense, but only in a “future” sense inasmuch as the children did not believe as yet, but were merely called to faith through the Sacrament. He thought that eventually they would attain to faith by virtue of the spiritual fellowship of children and sponsors, or rather the children and the Church; in his mind the child was brought to the baptismal font by the communion of saints, the Church, not by the sponsors only. During the period when the “scrutinies” enjoyed undisputed supremacy, the following order was observed in the baptism of infants. On the Monday after Oculi the parents were asked to announce the children to be baptized on Easter Sunday. On the following Wednesday the sponsors appeared in front of the church with the children. After their names were registered, the church doors were opened, and the sponsors moved in with their God-children lining up in two rows, the boys on the right side, the girls on the left side. Five acts made the children catechumens: the sign of the cross, laying on of hands, exorcism, presentation of salt, and prayer. The
\quellseite{77}
competent period was represented by seven scrutinies conducted in church in the presence of the children. The most important of these acts was the opening of the ears with which the \emph{traditio symboli et evangeliorum} and the \emph{redditio symboli} by the sponsor or the acolyte were connected. The last scrutiny was conducted on the Saturday before Easter. Later all these ceremonies were performed at one time so that the infants were brought to church but once; still, the fiction was maintained that in this one act the children passed through the two stages of the catechumenate and the photizomenate.

The training of children baptized in infancy and those who had remained unbaptized, was assigned by the Early Church exclusively to the Christian home. In this respect the example of Israel of old was followed where the home was expected to train the young. In the Jewish home reading was taught, the Ten Commandments, some of the Psalms, and important prayers were committed to memory (cf. Deut. 6:7); not rarely the Scriptures were read in the family circle. “From a child,” Paul writes to Timothy, “thou hast known the Scriptures” (2 Tim. 3:15). No change was made in this respect in the Jewish-Christian home, and the Gentile-Christian home soon acquired the same custom. Paul assumes that the children in Gentile-Christian congregations have learned the Ten Commandments (Eph. 6:1), and he urges the parents to bring their children up in the nurture and admonition of the Lord (Eph. 6:4; Col. 3:21). We know that for this purpose copies of the Bible were supplied to the Christian homes; and we have abundant proof for the wide prevalence of Bible reading in the Christian home, especially at noon or in the evening when the whole family was gathered around the table. Those who were too poor to buy copies were now and then visited by Christians who would read the Bible to them, also to the children. We
\quellseite{78}
hear of Origen that he “exercised already from childhood in the Scriptures” and that his father “made him each day learn and say by heart some passages” (Eusebius, \emph{Historia Eccles.} vi 2, 6 f., ed. by Schwartz, 1908). Something similar is told of Basil; and the \emph{Apost. Constit.} give the direction: “Teach your children thoroughly the word of the Lord and place in their hands every book of Holy Scripture” (iv 11). The singing and learning of Psalms was especially strongly recommended. Besides a number of Psalms, the children undoubtedly memorized the Creed and the Lord’s Prayer. The information thus acquired may have been somewhat fragmentary, yet this defect was largely supplied by the Christian attitude of the home.

Conditions changed materially after 325. Far and wide, parents began to neglect their sacred duties. It is this ominous situation which caused \emph{Chrysostom} who loved children and sympathetically appreciated their needs, so frequently to exhort to better care of the Church’s youth. No writer in the Early Church has discussed this problem so thoroughly. Without wearying and from ever new aspects, he appeals to the congregation to take serious this most urgent and essential of all tasks. He thinks that Christian parents will achieve the best results if they read the Scriptures together with their children, if they attend divine services with their children, and if they will ponder and contemplate the contents of the sermon with them. He desired the father to instruct his children in \emph{Biblical History}; and he did not think too highly of himself to write a booklet on the education of children. In this work he also offers practical examples of the most excellent type showing how a father should tell his children the stories of the Bible. Feeling his way, as it were, he laid down some of the most fundamental principles of teaching: the stories shall be selected according to the mental capacity of the children, they shall be presented
\quellseite{79}
in the form of free and vivid narratives with due attention to the limited range of the children’s concepts, in the teaching process the co-operation of the children shall be engaged, and the religious and moral truths of each story shall be set forth.

Chrysostom and others, as Basil, Rufinus, Salvian, also advised parents to have their children educated in \emph{convents}. Evidence is found that convents in the East and even more so in the West received children and gave them not only practical employment but also instruction. This applies chiefly, but by no means exclusively, to those who desired to become monks; many who later returned into the world to choose some secular calling were taught in convents. We lack more definite information about the religious instruction here given; Basil demands that instead of fables, the children should study “renowned events”, i.e., (to judge by the context) the events of Bible History, and that they should be made acquainted with the Proverbs of Solomon.

Only small numbers, however, could thus be reached. On the whole, indifference toward the most sacred duties advanced too swiftly to be checked. In proportion as truly Christian living decreased in the Church, fidelity in the instruction of the children vanished from the home. At the time when infant baptism became the general rule, the Church proved unequal to the task of devising new forms of education or of reforming the home. Thus the problem of religious instruction on which the future depended, was passed on unsolved to the medieval Church.
